\documentclass[11pt,a4paper]{article}
\usepackage[italian]{babel}
\usepackage[margin=2.5cm]{geometry}
\usepackage{parskip}
\usepackage{amsmath, amssymb, amsthm}
\usepackage{booktabs}
\usepackage{array}
\usepackage{xcolor}
\usepackage{needspace}
\usepackage{enumitem}
\usepackage{listings}
\usepackage{tikz}
\usepackage[most]{tcolorbox}
\usepackage[hidelinks]{hyperref}
\usetikzlibrary{decorations.pathreplacing, shapes.geometric, arrows.meta, positioning}

% Niente vedove/orfane: i paragrafi non lasciano righe isolate a fine/inizio pagina
\widowpenalty=10000
\clubpenalty=10000

% Elenchi più compatti
\setlist{itemsep=2pt, parsep=0pt, topsep=4pt, partopsep=0pt}

% --- Riquadri con bordo nero, senza numero (si spezzano tra due pagine) ---
% Uso: \begin{definizione} ... \end{definizione}
%      \begin{teorema}[Nome] ... \end{teorema}  (il nome è facoltativo)
\tcbset{nota/.style={enhanced jigsaw, breakable, colback=white,
  colframe=black, boxrule=0.5pt, sharp corners}}
\NewTColorBox{definizione}{}{nota,
  before upper={\textbf{Definizione.}\ }}
\NewTColorBox{teorema}{o}{nota,
  before upper={\textbf{Teorema\IfValueT{#1}{ (#1)}.}\ }}
\NewTColorBox{proposizione}{o}{nota,
  before upper={\textbf{Proposizione\IfValueT{#1}{ (#1)}.}\ }}
\NewTColorBox{proprieta}{o}{nota,
  before upper={\textbf{Proprietà\IfValueT{#1}{ (#1)}.}\ }}
\NewTColorBox{assioma}{o}{nota, boxrule=1pt,
  before upper={\textbf{Assioma\IfValueT{#1}{ (#1)}.}\ }}

% --- Ambienti semplici (senza riquadro) ---
% Il titolo sta in cima al blocco, anche se il contenuto inizia con un riquadro o
% con codice; e non resta mai da solo in fondo a una pagina.
% Uso: \begin{esempio}[Titolo facoltativo] ... \end{esempio}
\newcommand{\newrunin}[2]{%
  \NewDocumentEnvironment{#1}{o}{%
    \par\Needspace{4\baselineskip}\addvspace{\medskipamount}\noindent
    \textbf{#2}\IfValueT{##1}{ (##1)}\textbf{.}\ \ignorespaces}%
  {\par\addvspace{\medskipamount}}}
\newrunin{esempio}{Esempio}
\newrunin{esercizio}{Esercizio}
\NewTColorBox{osservazione}{}{nota,
  before upper={\textbf{Osservazione.}\ }}

% V e F nelle tavole di verità
\newcommand{\V}{\textsf{V}}
\newcommand{\F}{\textsf{F}}

% --- Codice C ---
% Uso: \begin{codice} ... \end{codice}      codice C numerato
%      \begin{comando} ... \end{comando}    comandi da terminale
%      \begin{risultato} ... \end{risultato} output di un programma
%      \lstinline|printf("ciao");|           codice dentro al testo
\lstdefinestyle{c}{
  language=C,
  basicstyle=\ttfamily\small,
  keywordstyle=\color{blue}\bfseries,
  commentstyle=\color{gray}\itshape,
  stringstyle=\color{red!70!black},
  numbers=left,
  numberstyle=\tiny\color{gray},
  numbersep=8pt,
  columns=flexible,
  keepspaces=true,
  breaklines=true,
  showstringspaces=false,
  tabsize=4,
  morekeywords={include, define},
  % lettere accentate nei commenti
  literate={à}{{\`a}}1 {è}{{\`e}}1 {é}{{\'e}}1 {ì}{{\`i}}1
           {ò}{{\`o}}1 {ù}{{\`u}}1 {È}{{\`E}}1 {À}{{\`A}}1
}
\lstset{style=c}
\newtcblisting{codice}{listing only, nota, left=6mm,
  listing options={style=c}}
\newtcblisting{comando}{listing only, nota,
  listing options={basicstyle=\ttfamily\small, numbers=none, language={},
    columns=flexible, keepspaces=true, breaklines=true}}
\newtcblisting{risultato}{listing only, enhanced jigsaw, breakable,
  colback=black!5, colframe=black, boxrule=0.5pt, sharp corners,
  listing options={basicstyle=\ttfamily\small, numbers=none, language={},
    columns=flexible, keepspaces=true, breaklines=true}}

% --- Diagrammi di flusso (simboli standard) ---
\tikzset{
  terminale/.style = {ellipse, draw, minimum width=2.5cm, minimum height=0.9cm},
  io/.style        = {trapezium, trapezium left angle=70, trapezium right angle=110,
                      draw, minimum width=2.5cm, minimum height=0.9cm},
  processo/.style  = {rectangle, draw, minimum width=2.5cm, minimum height=0.9cm},
  decisione/.style = {diamond, draw, aspect=2, minimum width=2.5cm},
  freccia/.style   = {thick, -{Stealth}}
}

\title{Esempi su cicli e istruzioni di selezione}
\author{Marco Cerbella}
\date{Lezione 17 -- 9 novembre 2026}

\begin{document}
\maketitle

\section{Perché i cicli sono importanti}

Un \emph{ciclo} è una delle strutture logiche di base della programmazione: definire un ciclo permette al calcolatore di eseguire lo stesso compito ripetutamente. Ogni volta che nella descrizione in linguaggio naturale delle specifiche del programma compaiono le parole «ogni» o «per ciascuno» serve un ciclo. Anche «per sempre» corrisponde a un ciclo: \lstinline|while(1)|.

\section{Esempi sui cicli}

\subsection{Condizioni logiche: \texttt{=} e \texttt{==}}

\begin{codice}
int x = 1;
while (x = 1) {          // ERRORE: assegnamento, non confronto
    //...
    printf("Insert 1 to continue");
    scanf("%d", &x);
}
\end{codice}
Se si usa un solo segno di uguale per controllare l'uguaglianza, il programma \emph{assegna} il valore a destra alla variabile a sinistra, e il risultato dell'espressione è il valore assegnato. In questo caso il valore è sempre $1$, che viene trattato come sempre vero: il ciclo non termina mai. La versione corretta è
\begin{codice}
int x = 1;
while (x == 1) {
    //...
    printf("Insert 1 to continue");
    scanf("%d", &x);
}
\end{codice}

\subsection{Trasformare i cicli}

Uno stesso ciclo si può scrivere con \lstinline|for|, \lstinline|while| o \lstinline|do...while|:
\begin{codice}
int n = 4;
for (int i = 0; i < n; i++) {
    printf("OK\n");
}
\end{codice}
\begin{codice}
int n = 4;
int i = 0;
while (i < n) {
    printf("OK\n");
    i++;
}
\end{codice}
\begin{codice}
int n = 4;
int i = 0;
do {
    printf("OK\n");
    i++;
} while (i < n);
\end{codice}

\subsection{do...while}

Il corpo di un \lstinline|do...while| viene eseguito \emph{almeno una volta}:
\begin{codice}
#include <stdio.h>
int main() {
    double number, sum = 0;
    // il corpo del ciclo viene eseguito almeno una volta
    do {
        printf("Enter a number: ");
        scanf("%lf", &number);
        sum += number;
    } while (number != 0.0);
    printf("Sum = %.2lf", sum);
    return 0;
}
\end{codice}

\subsection{Scorrere un array}

\begin{codice}
#include <stdio.h>
int main() {
    int foo[] = {16, 2, 77, 40, 12071};
    int n, result = 0;
    for (n = 0; n < 5; ++n) {
        result += foo[n];
    }
    printf("%d", result);      // stampa 12206
    return 0;
}
\end{codice}
La somma $16 + 2 + 77 + 40 + 12071$ dà $12206$.

\begin{osservazione}
Nelle slide l'esempio includeva \lstinline|<stdlib.h>| ma usava \lstinline|printf|; qui è stato corretto con \lstinline|<stdio.h>|.
\end{osservazione}

\subsection{Scorrere una matrice}

\begin{codice}
const int width = 5;
const int height = 3;
int main() {
    int mat[height][width];
    for (int n = 0; n < height; n++)
        for (int m = 0; m < width; m++)
            mat[n][m] = (n + 1) * (m + 1);
}
\end{codice}

\subsection{Array come argomenti di funzione}

Quando il nome di un array compare come argomento di una funzione, il compilatore lo converte implicitamente in un puntatore al primo elemento dell'array. Perciò \lstinline|int name[]| e \lstinline|int *name| sono equivalenti:
\begin{codice}
#include <stdio.h>
void fun(int a[], int n) {
    for (int i = 0; i < n; ++i)
        a[i] = i;
}
int main() {
    int a[3];
    fun(a, 3);
}
\end{codice}
\begin{codice}
#include <stdio.h>
void fun(int* a, int n) {
    for (int i = 0; i < n; ++i)
        *(a + i) = i;
}
int main() {
    int a[3];
    fun(a, 3);
}
\end{codice}
In C non esistono davvero «variabili array»: si lavora in realtà con puntatori, con una sintassi alternativa.

\subsection{Passare una matrice}

\begin{codice}
#include <stdio.h>
int matrix_sum(int row, int column, int m[row][column]) {
    int sum = 0;
    for (int i = 0; i < row; i++)
        for (int j = 0; j < column; j++)
            sum += m[i][j];
    return sum;
}

int main() {
    int m[2][3] = {1, 4, 5, 2, 1, 1};
    int result;
    result = matrix_sum(2, 3, m);
    printf("%d", result);      // stampa 14
}
\end{codice}
La matrice $2\times 3$ contiene per righe $\begin{smallmatrix}1&4&5\\2&1&1\end{smallmatrix}$, e la somma dei suoi elementi è $14$.

\subsection{Un esempio avanzato: cicli annidati (bubble sort)}

\begin{codice}
#include <stdio.h>
int main() {
    int n, c, d, swap;
    printf("Enter number of elements\n");
    scanf("%d", &n);
    printf("Enter %d integers\n", n);
    int array[n];
    for (c = 0; c < n; c++)
        scanf("%d", &array[c]);
    for (c = 0; c < (n - 1); c++)
    {
        for (d = 0; d < n - c - 1; d++)
        {
            if (array[d] > array[d+1])
            {
                swap       = array[d];
                array[d]   = array[d+1];
                array[d+1] = swap;
            }
        }
    }
    return 0;
}
\end{codice}
Il ciclo esterno esegue $n - 1$ passate; a ogni passata il ciclo interno confronta gli elementi adiacenti e li scambia se sono nell'ordine sbagliato, spostando così il massimo dei rimanenti in fondo all'array.

\section{Esempi sulle istruzioni di selezione}

\subsection{Condizioni logiche}

\begin{codice}
int age = 17;
if (age = 18)           // assegnamento!
    printf("OK!");
\end{codice}
Il valore dell'espressione \lstinline|age = 18| è $18$, che in C è vero: stampa sempre «OK!» e \lstinline|age| diventa $18$. Invece
\begin{codice}
int age = 17;
if (age == 18)
    printf("OK!");
\end{codice}
stampa «OK!» solo quando \lstinline|age| è uguale a $18$, e \lstinline|age| rimane $17$.

\subsection{Pari o dispari}

\begin{codice}
#include <stdio.h>
int main() {
    int number = 0;
    printf("Enter an integer: ");
    scanf("%d", &number);
    // vero se il resto è 0
    if (number % 2 == 0)
        printf("%d is an even integer.", number);
    else
        printf("%d is an odd integer.", number);
    return 0;
}
\end{codice}

\subsection{Catena di if}

Una sequenza di \lstinline|if| indipendenti \emph{non} equivale a una catena \lstinline|else if|. Nel primo caso tutte le condizioni vengono valutate; con \lstinline|score = 95| tutte sono vere e il voto finale sarà quello dell'ultimo \lstinline|if| soddisfatto:
\begin{codice}
char grade = 'F';
int score = 0;
scanf("%d", &score);
if (score >= 90) grade = 'A';
if (score >= 80) grade = 'B';
if (score >= 70) grade = 'C';
if (score >= 60) grade = 'D';
printf("Your grade is %c", grade);     // con 95 stampa D!
\end{codice}
Con la catena \lstinline|else if| solo il primo ramo vero viene eseguito:
\begin{codice}
char grade = 'D';
int score = 0;
scanf("%d", &score);
if (score >= 90)
    grade = 'A';
else if (score >= 80)
    grade = 'B';
else if (score >= 70)
    grade = 'C';
else if (score >= 60)
    grade = 'D';
printf("Your grade is %c", grade);
\end{codice}

\begin{osservazione}
Nelle slide l'ultimo confronto era scritto \lstinline|if (score >= 60);| con un punto e virgola finale: sarebbe un'istruzione vuota (errore comune descritto più sotto). Qui è stato tolto.
\end{osservazione}

\subsection{Attenzione all'annidamento}

Le due forme seguenti sono \emph{equivalenti}, perché l'\lstinline|else| si associa sempre all'ultimo \lstinline|if| senza \lstinline|else|:
\begin{codice}
if (x > 10) {
    if (x < 20)
        printf("x is between 10 and 20");
    else
        printf("x is greater than 20");
}
\end{codice}
\begin{codice}
if (x > 10)
    if (x < 20)
        printf("x is between 10 and 20");
    else
        printf("x is greater than 20");
\end{codice}
Invece, \emph{racchiudendo} il primo \lstinline|if| interno in un blocco, l'\lstinline|else| si associa all'\lstinline|if| esterno:
\begin{codice}
if (x > 10) {
    if (x < 20)
        printf("x is between 10 and 20");
}
else
    printf("x is less than 10");      // in realtà: x <= 10
\end{codice}

\subsection{Errori comuni}

\begin{enumerate}
    \item \emph{Punto e virgola dopo la condizione.} L'istruzione vuota diventa il corpo dell'\lstinline|if|, e la \lstinline|printf| viene eseguita sempre:
\begin{codice}
if (age >= 18);                       // il ';' chiude l'if
    printf("You can drink!\n");       // stampata sempre
\end{codice}
    \item \emph{Più istruzioni senza blocco.} Solo la prima istruzione dipende dall'\lstinline|if|; la seconda viene sempre eseguita:
\begin{codice}
if (age >= 18)
    printf("You can drink!\n");
    printf("Age greater than 18\n");  // stampata sempre
\end{codice}
    \item \emph{\lstinline|else| con condizione.} Un \lstinline|else| non può avere una condizione:
\begin{codice}
if (a < b)
    { /* ... */ }
else if (a > b)
    { /* ... */ }
else (a == b)                         // ERRORE: else non ha condizione
    { /* ... */ }
\end{codice}
    \item \emph{Confronti concatenati.} L'operatore \lstinline|>| è associativo da sinistra a destra, quindi con \lstinline|x = 15|:
\begin{codice}
int x = 15;
if (20 > x > 10)                      // (20 > x) > 10  ->  1 > 10  ->  falso
    printf("x is between 10 and 20");
\end{codice}
    L'espressione \lstinline|20 > x| è vera ($1$), e poi $1 > 10$ è falso ($0$). La forma corretta usa un operatore logico:
\begin{codice}
if (x > 10 && x < 20)
    printf("x is between 10 and 20");
\end{codice}
\end{enumerate}
Regole pratiche: (1) se ci sono più casi, un \lstinline|if| per ciascuno; (2) se la condizione è lunga, scomporla e usare gli operatori logici.

\subsection{Cicli infiniti e overflow}

\begin{codice}
#include <stdio.h>
#include <stdlib.h>
int main() {
    int i = 1;
    while (1) {
        i++;
        if (i < 0) {
            printf("%d\n", i);
            break;
        }
    }
}
\end{codice}
Il ciclo sembra infinito, ma \lstinline|i| continua a crescere fino a superare il massimo valore rappresentabile: la slide mostra l'output $-2147483648$. Questo è l'\emph{overflow}. Con un \lstinline|int| in complemento a due a $32$ bit:
\begin{center}
\small
\begin{tabular}{@{}ll@{}}
\texttt{01111111 11111111 11111111 11111111} & $2\,147\,483\,647$ \\
\texttt{01111111 11111111 11111111 11111110} & $2\,147\,483\,646$ \\
\dots & \dots \\
\texttt{00000000 00000000 00000000 00000001} & $1$ \\
\texttt{00000000 00000000 00000000 00000000} & $0$ \\
\texttt{10000000 00000000 00000000 00000000} & $-2\,147\,483\,648$ \\
\texttt{10000000 00000000 00000000 00000001} & $-2\,147\,483\,647$ \\
\dots & \dots \\
\texttt{11111111 11111111 11111111 11111111} & $-1$ \\
\end{tabular}
\end{center}
Aggiungendo $1$ al massimo $2\,147\,483\,647$ si passa al valore più negativo $-2\,147\,483\,648$.

\begin{osservazione}
Per gli interi \emph{con segno}, lo standard C considera l'overflow comportamento non definito: l'output mostrato è quello tipico su una macchina con complemento a due, non un risultato garantito.
\end{osservazione}

\end{document}
